\section{A sextic in thirty-five integers and an octic in thirty}
\label{sec:companions}

The same fixed coefficient tuple also gives two equations with fewer
unrestricted integer unknowns:
\[
 \begin{array}{c|r|r|r}
  \text{equation}&\text{unknowns}&\text{residuals}&\text{exact degree}\\\hline
  S_{35}&35&20&6\\
  O_{30}&30&15&8
 \end{array}
\]
Both satisfy the same ZFC-provable equivalence between absence of integer
zeros and $\operatorname{Con}(\mathrm{ZFC})$. The collection's historical
sextic and octic entries use forty-two and thirty-six unrestricted
integers, respectively. These are improvements of seven and six
unknowns relative to those entries, without a global minimality claim.

\subsection{Move a square guard to the coordinate that will be eliminated}

Let $Q_{42}^{X}$ have the same residuals as $Q_{42}$ except that
\[
 q-2\rho^2=0\qquad\text{is replaced by}\qquad X-8\xi^2=0.
\]
This replaces one unknown by one unknown. The retained equations include
$w=4u^2$ and $X=wq$.

\paragraph{Square-guard lemma.}
The integer zero sets of $Q_{42}$ and $Q_{42}^{X}$ are in bijection under
$\xi=u\rho$, with all other coordinates fixed.

\paragraph{Proof.}
The equations $F\alpha>0$ and $FC>0$ imply that $F,\alpha,C$ have the
same nonzero sign. Thus $q=F+\alpha+C$ satisfies $|q|\ge3$, before either
square guard has been used. If $u=0$, then $w=X=0$, the row
$E(X+1)=Xa$ gives $E=0$, and
\[
 a=(2v+1)q^3,\qquad |a|\ge27.
\]
With $H=4a+3$, the main norm has $D=ac+\gamma H$ and
$\Delta=a^2+4a+3\equiv a^2\pmod H$. Hence
$D^2-\Delta c^2=1$ implies $H\mid1$, contrary to $|H|\ge105$.
Therefore $u\ne0$.

At a zero of $Q_{42}^{X}$, the scale equations give
$qu^2=2\xi^2$. Write $\xi/u=A/B$ in lowest terms with $B>0$.
Then $qB^2=2A^2$, so $B^2\mid2$ and $B=1$. Consequently
$\rho=\xi/u$ is an integer and $q=2\rho^2$, restoring the deleted
guard. Conversely, a zero of $Q_{42}$ has
$X=4u^2\cdot2\rho^2=8(u\rho)^2$. The two substitutions are inverse.
\hfill$\square$

This elementary relocation matters for degree: $X=8\xi^2$ is quadratic,
while direct substitution of $w=4u^2$ into $X=wq$ would give a cubic
expression for $X$. The primary quartic remains the one in
Section~\ref{sec:quartic42}.

\subsection{The thirty-five-integer sextic, in full}

Keep the abbreviations $q=\beta J+1$, $C=q-F-\alpha$,
$k=\eta+\zeta$, and put $s=2v+1$. The twenty-six arithmetic unknowns are
\begin{gather*}
 J,F,\alpha,z,f,h,r,v,\tau_*,\eta,\zeta,y,\gamma,o,j,\\
 E,c,a,D,L,V,I,B_*,\Delta,u,\xi.
\end{gather*}
Together with the nine guard unknowns from Section~\ref{sec:quartic42},
there are thirty-five unknowns, all ranging over $\mathbb Z$.
Here $D,L,B_*,\Delta$ remain independent coordinates. The twenty
residuals are
\begin{align}
 S_1&=L-q^2+1, &
 S_2&=8\xi^2-4u^2q, \nonumber\\
 S_3&=a-E-sq(L+1), &
 S_4&=E(8\xi^2+1)-8\xi^2a, \nonumber\\
 S_5&=c-\eta-(a-E)k, &&\nonumber\\
 S_6&=\tau_*^2+(2\tau_*-k)E(c-\eta)-1, &&\nonumber\\
 S_7&=D-8\xi^2-ac-\gamma(4a+3), &&\nonumber\\
 S_8&=\Delta-a^2-4a-3, &
 S_9&=D^2-\Delta c^2-1.
 \label{eq:sextic35-first}
\end{align}
The other residuals are
\begin{align}
 S_{10}&=k-1-hE
    -\bigl[((q-1)(q-F)+\alpha)L+(M_C+qM_F)J\bigr], \nonumber\\
 S_{11}&=(K_0+4u^2)C-F-z(q-1), \nonumber\\
 S_{12}&=V+c-fo, &
 S_{13}&=V-jc+k-1-hE, \nonumber\\
 S_{14}&=I-(2r+1)c^2, &
 S_{15}&=B_*-I^2, \nonumber\\
 S_{16}&=B_*-\Delta(f^2-1), &
 S_{17}&=y^2-B_*(y^2-V^2)-1, \nonumber\\
 S_{18}&=F\alpha-\operatorname{Pos}(A_1,A_2,A_3), &&\nonumber\\
 S_{19}&=FC-\operatorname{Pos}(B_1,B_2,B_3), &&\nonumber\\
 S_{20}&=\eta\zeta-\operatorname{Pos}(T_1,T_2,T_3). &&
 \label{eq:sextic35-second}
\end{align}
The equation is
\begin{equation}
 \boxed{\displaystyle S_{35}=\sum_{i=1}^{20}S_i^2=0.}
 \label{eq:sextic35}
\end{equation}
Every residual has degree at most three. Only $S_{14}$ contains $r$,
and the coefficient of $r^2c^4$ in $S_{35}$ is exactly four. Thus its
degree is exactly six after the authentic coefficients are substituted.

To derive this system, first change the root coordinate by
\[
 \tau_{\mathrm{old}}=\tau_*+P-k,
\]
an invertible integral linear change. Then eliminate the seven registers
\[
 \begin{array}{c|l}
 P&k+E(c-\eta)\\
 F_2&f^2\\
 g&y^2-V^2\\
 X&8\xi^2\\
 N&q(L+1)\\
 C_2&c^2\\
 w&4u^2.
 \end{array}
\]
Each register has a defining equation with coefficient $1$ or $-1$.
After substitution, the old root restores as
$\tau_{\mathrm{old}}=\tau_*+E(c-\eta)$. The first norm becomes
\[
 [\tau_*+E(c-\eta)]^2
 -[k+E(c-\eta)]E(c-\eta)-1
 =\tau_*^2+(2\tau_*-k)E(c-\eta)-1.
\]
The cancellation of the fourth-degree term is exact. Every zero of the
twenty-residual system uniquely restores a zero of $Q_{42}^{X}$, and
every zero of $Q_{42}^{X}$ yields one in the other direction.

\subsection{The thirty-integer octic, in full}

The twenty-one arithmetic unknowns are now
\[
 J,F,\alpha,z,f,h,r,v,\tau,\eta,y,\gamma,o,j,E,c,a,V,I,u,\xi.
\]
Add the same nine guard unknowns. All thirty range over $\mathbb Z$.
The following are polynomial abbreviations, not independent unknowns:
\begin{gather*}
 q=\beta J+1,\quad C=q-F-\alpha,\quad s=2v+1,\quad H=4a+3,\\
 k=1-V+jc+hE,\quad b=E(c-\eta),\quad \Delta=a^2+4a+3,\\
 R=\bigl[(q-1)(q-F)+\alpha\bigr](q^2-1)+(M_C+qM_F)J.
\end{gather*}
The fifteen residuals are
\begin{align}
 O_1&=8\xi^2-4u^2q, &
 O_2&=a-E-sq^3, \nonumber\\
 O_3&=E(8\xi^2+1)-8\xi^2a, &
 O_4&=c-\eta-(a-E)k, \nonumber\\
 O_5&=\tau^2-(k+b)b-1, &&\nonumber\\
 O_6&=(8\xi^2+ac+\gamma H)^2-\Delta c^2-1, &&\nonumber\\
 O_7&=-V+jc-R, &
 O_8&=(K_0+4u^2)C-F-z(q-1), \nonumber\\
 O_9&=V+c-fo, &
 O_{10}&=I-(2r+1)c^2, \nonumber\\
 O_{11}&=I^2-\Delta(f^2-1), &
 O_{12}&=y^2-I^2(y^2-V^2)-1, \nonumber\\
 O_{13}&=F\alpha-\operatorname{Pos}(A_1,A_2,A_3), &&\nonumber\\
 O_{14}&=FC-\operatorname{Pos}(B_1,B_2,B_3), &&\nonumber\\
 O_{15}&=\eta(k-\eta)-\operatorname{Pos}(T_1,T_2,T_3). &&
 \label{eq:octic30-residuals}
\end{align}
The single equation is
\begin{equation}
 \boxed{\displaystyle O_{30}=\sum_{i=1}^{15}O_i^2=0.}
 \label{eq:octic30}
\end{equation}
Every residual has degree at most four. In $O_{11}$ the term
$-a^2f^2$ has coefficient $-1$. No other residual contains both $a$ and
$f$, so the coefficient of $a^4f^4$ in $O_{30}$ is exactly one. The
degree is therefore exactly eight on the fixed-coefficient slice.

Here the old root coordinate $\tau$ is retained. The twelve eliminated
coordinates restore, after composing the substitutions, as follows:
\[
 \begin{array}{c|l@{\qquad}c|l}
 N&q^3 & C_2&c^2\\
 F_2&f^2 & g&y^2-V^2\\
 w&4u^2 & X&8\xi^2\\
 \zeta&k-\eta & B_*&I^2\\
 D&8\xi^2+ac+\gamma H & \Delta&a^2+4a+3\\
 P&k+b & L&q^2-1.
 \end{array}
\]
The actual monic elimination order is
$N,C_2,F_2,g,w,X,\zeta,B_*,D,\Delta,P,L$.
In particular, $\zeta$ is first removed using
$\zeta=P-E(c-\eta)-\eta$, and $P$ is later removed using the exposed
congruence row. This yields the quadratic abbreviation for $k$, so the
common-sign guard still has degree only three.

Each elimination is an exact integral bijection on zero sets. The square
guard lemma then restores the original $Q_{42}$ system. Applying its
theorem proves the claimed consistency equivalences for both companions.

\paragraph{Reproducibility and scope.}
The companion generator checks nineteen monic eliminations, all fifty-four
parent-row restoration identities, all thirty-five factored residual
transcriptions, and both exact-degree coefficient certificates. It uses
explicit exceptions rather than removable Python assertions. Normal and
optimized runs agree. The literal and expanded JSON packets preserve the
whole coefficient binding; no new coefficient transfer is assumed.
These finite algebra checks supplement the ordinary proof. They do not
constitute a proof-assistant certificate for the inherited unbounded
Pell and compiler theorems. The JSON uses \texttt{rho} for $\xi$ and
\texttt{tau\_tail} for $\tau_*$.
